\medskip\noindent\textbf{命题 2.3.1。}\phantomsection \label{sga3:VIIB:2.3.1}{\itshape 设 \(K\) 和 \(G\) 为两个形式 \(k\)-群，其仿射代数分别为 \(B\) 和 \(A\)。假设 \(K\) 在 \(k\) 上拓扑平坦。则存在从 \[
  \operatorname{Hom}_{\operatorname{Grf}/k}(K,G)
\] 到满足下图交换的含幺 \(\mathbf O_k\)-代数同态 \[
  h:\mathbf H(K)\longrightarrow\mathbf H(G)
\] 的集合的典范双射：\begin{center}\phantomsection \label{sga3:VIIB:2.3.1:compatibility-diagram} % SGA3-VII-DIAG-102
% Authority: Exp7B-13oct24.pdf, local p. 48.
% Crop: figures/Exp7B_localp048_proposition231_compatibility_diagram.png
% Locator: tex/components/40_expose_VIIB_proposition231_covariant_algebra_morphisms_en.tex, Loop 1 line 32.
\begin{tikzcd}[row sep=1.25em,column sep=4em]
\mathbf H(K)\otimes_k\mathbf H(K)
  \arrow[r,"h\otimes h"]
  & \mathbf H(G)\otimes_k\mathbf H(G)
      \arrow[dr,"\varphi_G"] & {} \\
{} & {} & \mathbf H(G\times G) \\
\mathbf H(K) \arrow[uu,"\Delta_{\mathbf H(K)}"]
  \arrow[r,"h"]
  & \mathbf H(G) \arrow[ur,"\delta_G"']
  & {}
\end{tikzcd}
\end{center}}

由于 \(K\) 拓扑平坦，\(\mathbf H(K)\) 具有双代数结构 \textup{（参见 \SGARefLink{sga3:VIIB:2.2}{2.2}）}，且 \(\Delta_{\mathbf H(K)}\) 为其对角态射。换言之，用 \SGARefLink{sga3:VIIB:2.3}{2.3} 的记号，有 \[
  \Delta_{\mathbf H(K)}=\phi_K^{-1}\circ\delta_K.
\]。当 \(G\) 在 \(k\) 上也拓扑平坦时，命题由 \SGARefLink{sga3:VIIB:2.2.1}{2.2.1} 中建立的范畴等价推出。

在一般情形中，可以假设 \(k\) 为 Artin 环，并对代数 \[
  \mathbf H(K)=B^\dagger,
  \qquad
  \mathbf H(G)=A^\dagger.
\] 进行论证。以 \(\operatorname{Hom}_c(A,B)\) 表示从 \(A\) 到 \(B\) 的连续 \(k\)-线性映射的集合，以 \(\operatorname{Hom}_k(B^\dagger,A^\dagger)\) 表示从 \(B^\dagger\) 到 \(A^\dagger\) 的 \(k\)-线性映射的集合。\footnote{\label{sga3:VIIB:nde:87}编者补充了证明余下部分的细节。}

由 \SGARefLink{sga3:VIIB:0.3.6.A}{0.3.6.A}，若 \(M\)、\(P\) 为伪紧 \(k\)-模且 \(P\) 为投射模，则典范映射 \[
  \operatorname{Hom}_c(M,P)
  \longrightarrow
  \operatorname{Hom}_k(P^\dagger,M^\dagger),
  \qquad
  f\longmapsto{}^t f,
\] 为双射，其中 \({}^t f\) 表示 \(f\) 的转置。我们将分别在 \[
  M=A\widehat\otimes A,\quad P=B,
\] 或 \[
  M=A,\quad P=B\widehat\otimes B.
\] 的情形应用它。

设 \(f\in\operatorname{Hom}_c(A,B)\)。考虑下列图。标有 \textup{\SGARefLink{sga3:VIIB:0}{(0)}} 的方块交换，第二个图中两条未命名的竖箭头为 \({}^t(f\widehat\otimes f)\)。\begin{center}\phantomsection \label{sga3:VIIB:2.3.1:affine-diagram} % SGA3-VII-DIAG-103
% Authority: Exp7B-13oct24.pdf, local p. 48.
% Crop: figures/Exp7B_localp048_proposition231_affine_diagram.png
% Locator: tex/components/40_expose_VIIB_proposition231_covariant_algebra_morphisms_en.tex, Loop 1 line 87.
\begin{tikzcd}[row sep=3.1em,column sep=3.5em]
A\widehat\otimes A
  \arrow[r,"m_A"] \arrow[d,"{f\widehat\otimes f}"']
  \arrow[dr,phantom,"\textup{(1)}" description]
  & A \arrow[r,"\Delta_A"] \arrow[d,"f"]
      \arrow[dr,phantom,"\textup{(2)}" description]
  & A\widehat\otimes A \arrow[d,"{f\widehat\otimes f}"] \\
B\widehat\otimes B \arrow[r,"m_B"]
  & B \arrow[r,"\Delta_B"]
  & B\widehat\otimes B
\end{tikzcd}


  \medskip
  \phantomsection
  \label{sga3:VIIB:2.3.1:dual-diagram}
  % SGA3-VII-DIAG-104
  % SGA3-VII-DIAG-104
% Authority: Exp7B-13oct24.pdf, local p. 48.
% Crop: figures/Exp7B_localp048_proposition231_dual_diagram.png
% Locator: tex/components/40_expose_VIIB_proposition231_covariant_algebra_morphisms_en.tex, Loop 1 line 93.
\begin{tikzcd}[row sep=3.4em,column sep=2.15em]
A^\dagger\otimes A^\dagger
  \arrow[r,"\varphi_G"]
  \arrow[dr,phantom,"\textup{(0)}" description]
  & (A\widehat\otimes A)^\dagger
  & A^\dagger
      \arrow[l,"{\delta_G=m_A^t}"']
      \arrow[dr,phantom,"\textup{(2')}" description]
  & (A\widehat\otimes A)^\dagger
      \arrow[l,"\Delta_A^t"']
      \arrow[dr,phantom,"\textup{(0)}" description]
  & A^\dagger\otimes A^\dagger \arrow[l,"\varphi_G"'] \\
B^\dagger\otimes B^\dagger
  \arrow[r,equals] \arrow[u,"{{}^tf\otimes{}^tf}"]
  & (B\widehat\otimes B)^\dagger
      \arrow[u]
  & B^\dagger
      \arrow[l,"{\delta_K=m_B^t}"]
      \arrow[u,"{{}^tf}"']
      \arrow[ul,phantom,"\textup{(1')}" description]
  & (B\widehat\otimes B)^\dagger
      \arrow[l,"\Delta_B^t"]
      \arrow[u]
  & B^\dagger\otimes B^\dagger
      \arrow[l,equals] \arrow[u,"{{}^tf\otimes{}^tf}"']
\end{tikzcd}

\end{center}

若 \(f:A\to B\) 对应于形式群态射 \(K\to G\)，则方块 \textup{(1)} 和 \textup{(2)} 交换，且 \[
  \varepsilon_B\circ f=\varepsilon_A.
\]。因此方块 \textup{(1\(^{\prime}\))} 和 \textup{(2\(^{\prime}\))} 交换，且 \({}^t f\) 将 \(B^\dagger=\mathbf H(K)\) 的单位元映到 \(A^\dagger=\mathbf H(G)\) 的单位元。所以 \({}^t f\) 为含幺 \(k\)-代数同态 \[
  \mathbf H(K)\longrightarrow\mathbf H(G)
\]，且命题中的图交换。

反之，假设 \({}^t f\) 满足这些条件。则 \[
  \varepsilon_B\circ f=\varepsilon_A
\]，且方块 \textup{(1\(^{\prime}\))} 和 \textup{(2\(^{\prime}\))} 交换。对 \[
  M=A\widehat\otimes A,\quad P=B
  \qquad
  \text{（相应地，\(M=A,\ P=B\widehat\otimes B\)）},
\]，映射 \(g\mapsto{}^t g\) 为单射。因此方块 \textup{(1)} 和 \textup{(2)} 交换，故 \(f\) 与 \(A\) 和 \(B\) 的乘法及对角态射相容。

还须证明 \(f(1_A)=1_B\)。前面的论证给出 \[
  \varepsilon_B f(1)=1,
  \qquad
  \Delta_B f(1)=f(1)\widehat\otimes f(1),
  \qquad
  f(1)\cdot f(1)=f(1).
\]。由 \SGARefLink{sga3:VIIA:2.1}{2.1}\textup{(iii)}，前两个条件蕴含 \(f(1)\) 在 \(B\) 中的逆元为 \(c_Bf(1)\)。因此等式 \(f(1)\cdot f(1)=f(1)\) 蕴含 \(f(1)=1\)。所以 \[
  f:A\longrightarrow B
\] 是 \(\operatorname{Alp}/k\) 中的态射，且与 \(A\) 和 \(B\) 的余乘法相容。
